Подготовка научных и технических документов является неотъемлемой частью академической и профессиональной деятельности. Системы вёрстки документов позволяют авторам сосредоточиться на содержании, делегируя задачи форматирования специализированному программному обеспечению.

Система {\LaTeX}, созданная Лесли Лэмпортом на основе {\TeX} Дональда Кнута \cite{knuth1984texbook}, на протяжении десятилетий остаётся стандартом де-факто в академической среде~--- особенно в математике, физике и информатике \cite{latex-project}. Её мощные инструменты для набора математических формул и широкая экосистема пакетов делают её незаменимым инструментом для учёных и инженеров.

Вместе с тем {\LaTeX} обладает рядом существенных недостатков: сложный и многословный синтаксис, медленная компиляция для больших документов, нередко трудно интерпретируемые сообщения об ошибках. В 2023~году была представлена система Typst \cite{typst-github}~--- современная альтернатива, разработанная с учётом упомянутых ограничений.

В данной работе проводится сравнительный анализ систем {\LaTeX} и Typst по ключевым параметрам: синтаксис, производительность компиляции, набор математических формул и возможности экосистемы. Цель работы~--- оценить целесообразность перехода на Typst для типичных задач академической вёрстки.

Структура работы:
\begin{enumerate}
    \item в разделе~1 приводится обзор обеих систем;
    \item в разделе~2 выполняется сравнительный анализ по нескольким критериям;
    \item в разделе~3 рассматриваются практические примеры кода.
\end{enumerate}
